\documentclass[11pt]{article}
\input{packages.tex}
\DeclareMathOperator{\sign}{sign}
\DeclareMathOperator{\Imm}{Im}

% Настройка колонтитулов
\pagestyle{fancy}
\fancyhf{} % Очистить все поля
\fancyhead[L]{Какая гадость...какая гадость \sout{это ваше доказательство} эта ваша заливная рыба!}
\fancyhead[R]{\href{https://t.me/ceagest}{Автор}}
\fancyfoot[C]{\thepage}               % Справа внизу

\setlength{\headheight}{13.6pt}

\hypersetup{
    colorlinks=true,
    linkcolor=blue,
    citecolor=blue,
    filecolor=blue,
    urlcolor=blue,
}

\begin{document}

\begin{titlepage}
    \centering
    \vspace*{\fill}  % заполняет пространство сверху

    {\Huge Лекция 3. Исчисление высказываний и исчисление предикатов. Их синтаксис и семантика. Выводимость.
    \par}

    \vspace*{\fill}  % заполняет пространство снизу
\end{titlepage}

\setcounter{page}{1}

\section{Исчисление высказываний}

\subsection{Синтаксис исчисления высказываний. Выводимость}

\paragraph{Алфавит.} 
\begin{enumerate}
    \item Счётное множество символов переменных
    \item «$\lnot$», «$\to$», «$($», «$)$»
\end{enumerate}

\paragraph{Формулы.}
\begin{itemize}
    \item Если $A$ -- символ переменной, то $A$ -- формула.
    \item Если $A = (\lnot B)$, где $B$ -- формула, то $A$ -- формула.
    \item Если $A = (B \to C)$, где $B$ и $C$ -- формулы, то $A$ -- формула.
    \item Других формул нет (не всегда уточняют, но мы это делаем).
\end{itemize}

\paragraph{Аксиомы.} 
Аксиом бесконечно много, поэтому зададим их схемами:
\begin{enumerate}
    \item $(A \to (B \to A))$
    \item $((A \to (B \to C)) \to ((A \to B) \to (A \to C)))$
    \item $(((\lnot B) \to (\lnot A)) \to (((\lnot B) \to A) \to B))$
\end{enumerate}
Вместо $A, B$ и $C$ можно подставлять любые формулы, но для одинаковых букв 
должны быть одинаковые формулы.

\paragraph{Правила вывода.} 
\underline{\textit{modus ponens}}: $A, (A \to B) \vdash B$

\begin{definition}[Выводимость]
    Говорят, что формула $A$ \textbf{выводится} в исчислении высказываний, если 
    существует последовательность $A_1, A_2, \ldots, A_n = A$ такая, что каждое $A_i$ либо является аксиомой, либо выводится из предыдущих посредством одного применения modus ponens, то есть 
    \[\exists j, k < i: A_k = A_j \to A_i\]
    Число $n$ будем называть \textbf{длиной вывода}.
\end{definition}

\begin{remark}
    Если в определении выше заменить словосочетания "выводится в исчислении высказываний" $\;$ и "modus ponens" $\;$ на "выводится в формальной системе" $\;$ и "некоторого правила вывода" $\;$ соответственно, 
    то получим определение выводимости в произвольной формальной системе. Мы воспользуемся этим при формулировке выводимости в исчислении предикатов.
\end{remark}

\begin{definition}[Условная выводимость]
    Пусть $\Gamma$ --- множество формул. Будем говорить, что формула $A$ выводится из $\Gamma$, если 
    существует последовательность $A_1, A_2, \ldots, A_n = A$ такая, что каждое $A_i$ либо является аксиомой, либо является гипотезой, либо выводится из предыдущих посредством одного применения modus ponens, то есть 
    \[\exists j, k < i: A_k = A_j \to A_i\]
\end{definition}

\begin{example}
    Давайте выведем формулу $(A \to A)$. Для этого, согласно определению, построим нужную последовательность, используя схемы аксиом и modus ponens:
    \begin{enumerate}
        \item \underline{\textnormal{A1}} \textnormal{:} $A \to ((A \to A) \to A)$
        \item \underline{\textnormal{A2}} \textnormal{:} $(A \to ((A \to A) \to A)) \to ((A \to (A \to A)) \to (A \to A))$
        \item \underline{\textnormal{MP 1,2}} \textnormal{:} $(A \to (A \to A)) \to (A \to A)$
        \item \underline{\textnormal{A1}} \textnormal{:} $A \to (A \to A)$
        \item \underline{\textnormal{MP 3,4}} \textnormal{:} $A \to A$
    \end{enumerate}
\end{example}

\begin{example}
    Выведем теперь из формулы $A$ формулу $B \to A$, пользуясь определением условной выводимости (то есть здесь полагаем $\Gamma = \{A\}$):
    \begin{enumerate}
        \item \underline{\textnormal{A1}} \textnormal{:} $A \to (B \to A)$
        \item $A$ --- гипотеза
        \item \underline{\textnormal{MP 1,2}} \textnormal{:} $B \to A$
    \end{enumerate}
\end{example}

\subsection{Семантика исчисления высказываний}

Естественной семантикой для исчисления высказываний является \textbf{булева алгебра}: 
переменным сопоставляется значение из множества $\{0, 1\}$.

\begin{definition}
    \textbf{Тавтологией} будем называть формулу, принимающую значение $1$ при любом наборе значений переменных.
\end{definition}

\begin{theorem}[Корректность исчисления высказываний]
    В исчислении высказываний выводимы (без использования гипотез) только тавтологии.
\end{theorem}

\begin{theorem}[Полнота исчисления высказываний]
    В исчислении высказываний выводимы \textbf{все} тавтологии.
\end{theorem}

\begin{definition}
    Будем говорить, что формула $A$ \textbf{семантически} следует из множества формул $\Gamma$, и записывать это $\Gamma \vDash A$, если $A$ истинна 
    на всех наборах значений переменных, на которых истинны \textbf{все} формулы из множества $\Gamma$.
\end{definition}

\begin{remark}
    Вообще говоря, на всех тех наборах, для которых истинны не все формулы из $\Gamma$, формула может быть как истинна, так и ложна.
\end{remark}

\begin{theorem}[Критерий условной выводимости]
    Пусть $\Gamma$ --- множество формул, $A$ --- формула.
    \[\Gamma \vdash A \iff \Gamma \vDash A\]
\end{theorem}

\begin{theorem}[О дедукции]
    Пусть $\Gamma$ --- множество формул, $A$ и $B$ --- формулы.
    \[\Gamma \cup \{A\} \vdash B \iff \Gamma \vdash A \to B\]
\end{theorem}

\section{Исчисление предикатов}

\subsection{Синтаксис исчисления предикатов. Выводимость}

\paragraph{Алфавит.} 
\begin{enumerate}
    \item Счётное множество символов переменных $x$
    \item Счётное множество функциональных символов $\mathcal{F}$
    \item Счётное множество символов предикатов $\mathcal{P}$
    \item «$\lnot$», «$\to$», «$($», «$)$», «$\forall$», «$,$»
\end{enumerate}

\paragraph{Формулы.}

\begin{definition}
    Определим \textbf{терм} рекурсивно:
    \begin{itemize}
        \item Если $t$ --- символ переменной, то $t$ --- терм
        \item Если $t_1, \ldots, t_k$ --- термы и $f \in \mathcal{F}_k$ (то есть $f$ является $k$-арным (имеющим $k$ аргументов) функциональным символом), то $f(t_1, \ldots, t_k)$ ---  терм
    \end{itemize}
\end{definition}

\begin{definition}
    \textbf{Атомарной формулой} будем называть $A(t_1, \ldots, t_k)$, где $A \in \mathcal{P}_k$ (то есть $A$ является $k$-арным символом предиката), $t_1, \ldots, t_k$ --- термы. 
\end{definition}

\begin{definition}
    Определим \textbf{формулу} рекурсивно:
    \begin{itemize}
        \item Если $A$ --- атомарная формула, то $A$ --- формула
        \item Если $A = (\lnot B)$, где $B$ --- формула, то $A$ --- формула
        \item Если $A = (B \to C)$, где $B$ и $C$ --- формулы, то $A$ --- формула
        \item Если $A = \forall x B$, где $x$ --- символ переменной, $B$ --- формула, то $A$ --- формула
        \item Других формул нет
    \end{itemize}
\end{definition}

\begin{definition}
    Пусть $x$ --- символ переменной, $A$ --- формула, $(\forall x A)$. Тогда формула $A$ называется \textbf{областью действия} квантора по переменной $x$.
\end{definition}

\begin{definition}
    Вхождение переменной в формулу называется \textbf{свободным}, если оно не лежит в области действия квантора по данной переменной. Если 
    вхождение переменной не свободно, то его называют \textbf{связанным}.
\end{definition}

\begin{example}
    Пусть $x, y$ --- символы переменных, $T \in \mathcal{P}_2$, $f \in \mathcal{F}_2$. Тогда все вхождения переменной $x$ в формулу $\forall x T(f(x, y), x)$ являются связанными, а все вхождения $y$ --- свободные.
\end{example}

\begin{remark}
    В формулах, подобных формуле $\displaystyle \mathop{\underline{\forall x}}_{(1)} A \to \mathop{\underline{\forall x}}_{(2)} B$, можно мыслить $(1)$ и $(2)$ как программу, в которой идут два \textbf{не связанных между собой} счётчика по $x$.
\end{remark}

\paragraph{Аксиомы.} 
Аксиом бесконечно много, поэтому зададим их схемами:
\begin{enumerate}
    \item $(A \to (B \to A))$
    \item $((A \to (B \to C)) \to ((A \to B) \to (A \to C)))$
    \item $(((\lnot B) \to (\lnot A)) \to (((\lnot B) \to A) \to B))$
    \item $\forall x (A \to B) \to (A \to \forall x B)$, при условии, что в $A$ нет свободных вхождений $x$
    \item $\forall x A \to A[x := t]$, при условии, что ни одно свободное вхождение $x$ в $A$ не лежит в области действия кванторов по переменным, входящим в $t$.
\end{enumerate}
Здесь $A, B$, $C$, $t$ --- некоторые формулы, $x$ --- символ переменной. Запись $A[x := t]$ означает результат подстановки $t$ в $A$ вместо всех свободных вхождений $x$. 

\paragraph{Правила вывода.} 
\begin{enumerate}
    \item \underline{\textit{modus ponens}}: $A, (A \to B) \vdash B$
    \item \underline{\textit{generalization}}: $A \vdash \forall x A$
\end{enumerate}

\begin{definition}[Выводимость]
    Говорят, что формула $A$ \textbf{выводится} в исчислении предикатов, если 
    существует последовательность $A_1, A_2, \ldots, A_n = A$ такая, что каждое $A_i$ либо является аксиомой, либо выводится из предыдущих посредством одного применения некоторого правила вывода в исчислении предикатов.
    Число $n$ будем называть \textbf{длиной вывода}.
\end{definition}

\begin{definition}[Условная выводимость]
    Пусть $\Gamma$ --- множество формул. Будем говорить, что формула $A$ выводится из $\Gamma$, если 
    существует последовательность $A_1, A_2, \ldots, A_n = A$ такая, что каждое $A_i$ либо является аксиомой, либо является гипотезой, либо выводится из предыдущих посредством одного применения некоторого правила вывода в исчислении предикатов.
\end{definition}

\subsection{Семантика исчисления предикатов}

\begin{definition}
    \textbf{Сигнатура} --- набор предикатных или функциональных символов с указанием арности.
\end{definition}

\paragraph{Интерпретация сигнатуры.} Пусть $\mathcal{M}$ --- область интерпретации (проще говоря, множество, из которого берём значения переменных). Каждому функциональному символу сопоставим функцию из $\mathcal{M}^k$ в $\mathcal{M}$.
Каждому предикатному символу $T \in \mathcal{P}_k$ сопоставим $T_k \subseteq \mathcal{M}^k$.

\begin{definition}
    Формула называется \textbf{общезначимой}, если она истинна на любой интерпретации при любом наборе значений переменных.
\end{definition}

\end{document}